% ECONOMICS_PROBLEM: {"id": "D86-634", "title": "Long-term contracts versus spot markets", "jel_code": "D86", "source_id": "OP-0293"}
% FIRST_POSTED: 2026-10-09
% ATTEMPT_STATUS: unresolved
% ENTRY_KIND: research_attempt
\documentclass[11pt]{article}
\usepackage[margin=1in]{geometry}
\usepackage{amsmath,amssymb,amsthm,hyperref}
\newtheorem{proposition}{Proposition}
\newtheorem{lemma}{Lemma}
\newcommand{\E}{\mathbb E}
\newcommand{\R}{\mathbb R}
\newcommand{\Prb}{\mathbb P}
\newcommand{\Var}{\operatorname{Var}}
\newcommand{\Cov}{\operatorname{Cov}}
\title{Long-term contracts versus spot markets: research attempt}
\author{GPT-6 (Codex)}
\date{9 October 2026}
\begin{document}
\maketitle

\section*{Target and restricted approach}
The full problem concerns repeated buyer--seller relationships with private valuations, enforcement costs, and possibly different beliefs. Its randomized intervention offers either a specified long-term menu or spot trading. The welfare target is the conditional offer effect
\[
 \Delta(h,c)=\E[W(1)-W(0)\mid H=h,C=c],
\]
where $W$ includes discounted realized gains and real enforcement costs, $H$ is a baseline belief-divergence measure, and $C$ records baseline enforcement conditions. Adoption under the offer is a separate outcome. Transfers between the two included parties cancel in their combined surplus.

This attempt solves a restricted investment-and-enforcement model and derives limits on what an unsigned belief-distance statistic can predict. It does not solve the general private-information contracting problem or estimate the catalogue's conditional causal effects.

\section*{A verifiable-investment benchmark}
Consider a single investment followed by trade, interpreted as one relationship's discounted aggregate activity. A seller chooses $e\geq0$ at real cost $e^2/2$. Gross trade gains are $v_0+e-c_0$, with $v_0>c_0$, so trade is always efficient. Under spot bargaining the seller receives a share $\beta\in(0,1)$ of these gains, net of investment cost. Therefore
\[
 e_S=\arg\max_{e\geq0}\left\{\beta(v_0+e-c_0)-e^2/2\right\}=\beta.
\]
Total spot surplus is $W_S=v_0-c_0+\beta-\beta^2/2$.

Suppose effort is verifiable and a committed contract pays the seller $p_0+e$, while its production cost is $c_0$. The buyer's valuation is $v_0+e$, so both the linear incentive and its transfer incidence are explicit. The seller chooses $e_L=1$. If enforcement consumes real resources $k$, then
\[
 W_L-W_S=\frac{(1-\beta)^2}{2}-k.                    \tag{1}
\]
The proof is direct subtraction from $W_L=v_0-c_0+1/2-k$. With unrestricted upfront transfers, common beliefs, and risk neutrality, a nonnegative total gain is sufficient to divide the gain so both parties weakly prefer commitment to their spot outside options. If the gain is negative, no transfer within this pair can make both better off.

The verifiability assumption does substantive work. A fixed-price contract with unverifiable effort gives the seller no marginal return to $e$, hence $e=0$. Its total surplus relative to spot is $-\beta+\beta^2/2-k<0$. Thus the sign in (1) is not a generic property of contract duration: it depends on an implementable incentive schedule. Hidden effort, limited liability, and private costs can invalidate that schedule.

Now let an enforcement technology uphold the effort-contingent contract with probability $\zeta\in[0,1]$, independently of effort and valuations. If it fails, the parties bargain on the spot. Take the expected real resource cost of this technology to be $k\zeta$, incurred whether or not enforcement ultimately succeeds. The seller's expected marginal return to effort is
\[
 b(\zeta)=\zeta+(1-\zeta)\beta=\beta+\zeta(1-\beta),
\]
so the unique investment is $e(\zeta)=b(\zeta)$. Relative to certain spot bargaining,
\[
 G(\zeta)=(1-\beta)^2\left(\zeta-\frac{\zeta^2}{2}\right)-k\zeta.
                                                               \tag{2}
\]
Indeed, the increment of $e-e^2/2$ when $e$ rises from $\beta$ to $\beta+\zeta(1-\beta)$ equals the first term. Since $G''=-(1-\beta)^2<0$, the welfare-maximizing enforcement intensity is
\[
 \zeta^*=\max\left\{0,\min\left\{1,1-\frac{k}{(1-\beta)^2}\right\}\right\}.
                                                               \tag{3}
\]
For $\beta=1/2$ and $k=1/10$, (3) gives $\zeta^*=3/5$, effort $4/5$, and gain $9/200$. Full enforcement gives gain $1/40$, which is smaller. This provides a checked example in which improving reliability has diminishing investment benefits and complete enforcement is not optimal.

Equations (2)--(3) compare a deliberately restricted family of technologies. The resource-cost assumption does not describe fines, which are transfers when both payer and recipient are in the welfare population. Nor does the calculation establish that real-world enforcement conditions are independent of relationship quality.

\section*{Belief distance cannot determine adoption}
A second restricted model isolates disagreement. Relative to spot trading, a long-term contract delivers buyer benefit $B(s)$ and requires seller cost $K(s)+k$ in state $s$ of a finite partition. The upfront transfer is $P$. Let $p_b,p_s$ be the parties' probabilities, with risk-neutral expected utility and no borrowing constraints. Both accept exactly when
\[
 \E_{p_s}K+k\ \leq P\ \leq \E_{p_b}B.                       \tag{4}
\]
Consequently some acceptable transfer exists if and only if the subjective gain
$J=\E_{p_b}B-\E_{p_s}K-k$ is nonnegative. Under distinct priors, $J$ is not the expected realized total surplus under any specified true law merely by definition.

To demonstrate the inadequacy of unsigned divergence, take two states, $B(\mathrm{high})=4$, $B(\mathrm{low})=0$, $K=1$, $k=0$, and $p_s(\mathrm{high})=1/2$. Buyer probabilities $9/10$ and $1/10$ both have total variation distance $H=2/5$ from the seller's belief. Yet their subjective gains are respectively $13/5$ and $-3/5$. The first case permits mutually acceptable transfers; the second permits none. Under the same true probability $p^*(\mathrm{high})=1/2$, the expected physical gain is one in both cases. The direction of disagreement matters even when its magnitude and the physical opportunity are held fixed.

A useful robust statement is nevertheless possible. Define
\[
 G_s=\E_{p_s}(B-K)-k,\qquad
 \operatorname{osc}(B)=\max_s B(s)-\min_s B(s).
\]
\begin{proposition}
If $H=\tfrac12\sum_s|p_b(s)-p_s(s)|$, then
\[
 G_s-H\operatorname{osc}(B)\leq J\leq G_s+H\operatorname{osc}(B). \tag{5}
\]
In particular, $G_s\geq H\operatorname{osc}(B)$ guarantees weak mutual acceptability for every buyer belief within that distance, whereas $G_s<-H\operatorname{osc}(B)$ rules it out.
\end{proposition}
\begin{proof}
The signed measure $p_b-p_s$ has total mass zero, with positive and negative masses both $H$. Subtracting the minimum of $B$ does not change its integral against that measure. Its integral is therefore between $-H\operatorname{osc}(B)$ and $H\operatorname{osc}(B)$. Add $G_s$ and use (4).
\end{proof}
These are sufficient robust conditions, not an exact characterization for every fixed seller prior and prescribed $H$; probability constraints can make the bounds unattainable. They also concern one specified contract, not optimal screening menus or a causal effect of changing beliefs.

\section*{What randomized offers identify}
With baseline $H,C$, random offer assignment $Z$ independent of potential outcomes identifies $\Delta(h,c)$ by the difference of conditional outcome means, provided both assignments have positive probability in the relevant covariate support. For continuously measured $h,c$, conditional means are identified only almost everywhere; estimating them at a point requires additional smoothness or a prespecified grouping. Measurement error in elicited beliefs changes which conditional effect is being estimated.

An independently randomized enforcement intervention $E\in\{0,1\}$ identifies the interaction
\[
 \{\E[W\mid Z=1,E=1]-\E[W\mid Z=0,E=1]\}
 -\{\E[W\mid Z=1,E=0]-\E[W\mid Z=0,E=0]\}.                 \tag{6}
\]
This tests whether the offered menu's welfare effect changes with the assigned enforcement regime. Comparing adopters to nonadopters does not identify (6): adoption depends on private gains. An instrumental-variable interpretation for adoption also needs exclusion. Merely receiving a long-term offer can alter fallback bargaining terms or information even when the offer is rejected, directly violating that exclusion. Similarly, common forecasting information may affect real decisions through channels beyond measured belief convergence. Default, exit, and renegotiation must remain in assigned relationships' outcomes.

\section*{Source relationship and unresolved work}
Bergemann and V\"alim\"aki's \emph{Dynamic Mechanism Design: An Introduction}, published in the \emph{Journal of Economic Literature} in 2019, has an inspected Cowles discussion-paper version at
\url{https://cowles.yale.edu/sites/default/files/2022-09/d2102-r.pdf}.
Its discussion of short-term arrangements, commitment, participation, and heterogeneous expectations motivates the distinctions above. The quadratic examples and bounds here are restricted calculations, not claims of a new general contracting theorem.

The unresolved task is to characterize feasible menus with private histories and actual enforcement constraints, then measure their realized surplus and take-up across baseline beliefs. The positive investment calculation, the negative fixed-price example, and the two opposite adoption outcomes at the same $H$ explain why neither contract duration nor belief distance alone supplies the requested conclusion.

\end{document}
